% =============================================================
%  hydrothermal_goethite_abstract.tex
%  Abstract for the merged (Ch1 + Ch2) GCA manuscript.
%  Kept separate from any \begin{abstract} wrapper, same convention
%  as chapter1_abstract.tex and chapter2_abstract.tex.
%
%  Working title:
%    High Helium Retentivity in Goethite under Hydrothermal
%    Conditions, Revealed by Step-Heat Experiments in Silica Glass
%    Ampoules
%
%  Note the ordering: result first, method second. GCA does not
%  publish method development unless it carries an enabled scientific
%  result, so the goethite finding must lead.
% =============================================================

U-Th/He geochronology of goethite (FeOOH) dates weathering and ferruginization directly, 
but interpreting a goethite date requires knowing how the mineral retains radiogenic helium, 
which has proven difficult to measure. Recent work revealed that goethite decomposes to hematite 
and water during vacuum step heating, and helium release records dehydration kinetics rather than 
helium migration through stable goethite. We report the first helium step-heat experiment on goethite 
under hydrothermal conditions, made possible by sealing the sample and liquid water inside a pure 
silica glass ampoule, which retains water at up to 85\,bar while transmitting helium to the vacuum system. 
Permeation through the ampoule wall delays measurement of released helium; we quantify this lag with a 
first-order permeation model and remove it either by measuring the ampoule's residual helium with low-temperature 
permeation steps or, for long steps, by bounding it. Encapsulated Durango apatite reproduces its published Arrhenius array 
under this correction, establishing that the ampoule contribution can be separated from the sample's. 
The ampoule's permeation rate constant is measured during the experiment itself from series of isothermal
permeation steps, which is necessary because water-bearing ampoules permeate faster than published rates for
anhydrous silica glass rather than more slowly; we attribute this to displacement of gas-phase volume by liquid
water rather than to any change in the glass, and it makes published permeabilities a conservative lower bound.
Proton-irradiated goethite Win-06-01b, previously step heated in vacuum, was
held at water saturation over 170--250\textdegree C for 40 days through multiple prograde and retrograde cycles. 
Apparent helium diffusivities are approximately seven natural log units below those measured on the same sample 
in vacuum. Above 7\% cumulative \textsuperscript{3}He yield the array is linear, 
with an activation energy of \ChTwoEa\ $\pm$ \ChTwoEaErr\,kJ/mol, and its delta plot follows 
multiple diffusion domain systematics. A domain with these parameters is nearly quantitatively retentive of helium at 25\textdegree C over billion-year timescales. Because the difference between hydrothermal and vacuum experiments on the same aliquot exceeds the differences among vacuum experiments on unrelated samples, vacuum step heats of goethite appear to record the mineral's response to the experiment rather than helium migration through it, and helium retention in natural goethite is correspondingly more robust than those experiments imply. The method extends directly to systematic tests of crystallite size, chemical composition, and radiation damage.
